% !TEX root = ../main.tex

%%% ============================================================
\section{Quá trình tạo và hoàn thiện báo cáo}
%%% ============================================================

\subsection{Năm giai đoạn viết báo cáo}
\begin{frame}{Báo cáo được viết như thế nào: 5 giai đoạn}
  \footnotesize
  Nhóm dùng hai công cụ AI với hai vai trò tách biệt: \textbf{một công cụ để viết}
  (brainstorm dàn ý, viết nháp, sửa file, build PDF) và \textbf{một công cụ khác để
  đọc lại và phản biện} bản vừa viết ra --- giống như nhờ một người khác đọc bản
  thảo thay vì tự mình vừa viết vừa tự chấm điểm cho chính mình. \textbf{Sinh viên
  là người quyết định} ở mọi bước: chọn giữ ý nào, bỏ ý nào, viết lại câu chữ cuối
  cùng bằng văn phong của mình.

  \vspace{0.15cm}
  \textbf{Ví dụ cụ thể — viết Chương 5 (Bàn luận)}:
  \begin{enumerate}\scriptsize\setlength{\itemsep}{2pt}
    \item \textbf{Dàn ý}: đưa cho công cụ viết bốn câu hỏi nghiên cứu (CH1--CH4) và
      số liệu thực nghiệm đã có $\rightarrow$ nó đề xuất dàn ý 3 mục; sinh viên chốt lại.
    \item \textbf{Viết nháp}: công cụ viết soạn từng mục theo dàn ý đã chốt, chỉ dùng
      số liệu có thật trong file kết quả, không được tự suy diễn con số.
    \item \textbf{Đọc lại phản biện}: công cụ đọc lại phát hiện mục CH4 chưa thực sự trả
      lời câu hỏi đặt ra ở Chương 1 $\rightarrow$ chỉ ra chỗ thiếu, đề xuất hướng bổ sung.
    \item \textbf{Sinh viên chỉnh sửa}: đọc góp ý, quyết định bổ sung đoạn nào, viết lại
      bằng lời văn của mình, rồi mới đưa công cụ viết sửa file \texttt{.tex}.
    \item \textbf{Build lại \& lưu}: build PDF kiểm tra không lỗi, rồi mới commit.
  \end{enumerate}
\end{frame}

\subsection{Vòng chỉnh sửa theo nhận xét của GVHD}
\begin{frame}{Một vòng chỉnh sửa diễn ra như thế nào}
  \begin{columns}[t]
    \begin{column}{0.52\textwidth}
      \footnotesize
      \textbf{Các bước, minh họa bằng một lần sửa thật}
      \begin{enumerate}\scriptsize\setlength{\itemsep}{2.5pt}
        \item GVHD nhận xét: "kết quả và bàn luận đang lẫn vào nhau, tách riêng ra".
        \item Đọc bản hiện tại cùng nhận xét đó, liệt kê những chỗ cụ thể cần sửa
          và xếp thứ tự quan trọng.
        \item Sửa từng chỗ trong file \texttt{.tex}, build lại PDF sau mỗi lần sửa
          để kiểm tra ngay, không dồn hết vào cuối.
        \item Đọc lại lần hai: \textbf{xác nhận trước} chỗ nào đã ổn, rồi mới nêu
          thêm điều còn thiếu, chia theo mức độ ưu tiên.
        \item Còn điểm cần sửa thì quay lại bước 3; khi ổn thì gửi lại cho GVHD.
      \end{enumerate}
    \end{column}
    \begin{column}{0.44\textwidth}
      \begin{alertblock}{Những gì luôn được kiểm tra lại}
        \footnotesize
        \begin{itemize}\setlength{\itemsep}{2.5pt}
          \item \textbf{Trích dẫn}: mỗi nguồn được tra lại trên Google Scholar/DOI;
            không tìm thấy, hoặc bài báo không nói đúng điều đang trích, thì bỏ hoặc sửa lại.
          \item \textbf{Số liệu}: mọi con số trong báo cáo chỉ lấy từ kết quả chạy
            thực tế, không được bịa hay chỉnh sửa cho đẹp.
          \item \textbf{Vai trò của AI}: AI chỉ là công cụ hỗ trợ, không phải đồng
            tác giả; sinh viên chịu trách nhiệm toàn bộ nội dung.
        \end{itemize}
      \end{alertblock}
      {\scriptsize Trước mỗi lần nộp: build lại toàn bộ, đảm bảo không còn lỗi biên dịch hay tham chiếu thiếu.}
    \end{column}
  \end{columns}
\end{frame}
